\section{Only two channels need to cross the extra learner layers}
\label{sec:two-channel-padding}
The last hidden width of a comparison implementation can be much
smaller than its largest hidden width. Scalar output permits a useful
compiler: after completing the teacher computation, carry its positive
and negative parts through the additional layers.

\begin{lemma}[Two-channel output compiler]
\label{lem:two-channel-compiler}
Let a ReLU teacher have hidden depth $h$, width at most $m$, linear
scalar readout $z_g$, coefficient cap $B\ge1$, and rooted formula
size $S$ before the readout. There is an implementation of the same
raw function with hidden depth $h+1$, width at most $m$,
coefficient cap $B$, and rooted formula size at most $2S+2$.
Its last hidden layer has at most two relevant neurons and its final
nonzero final readout coefficients belong to $\{-1,1\}$.
At any larger hidden depth $H$, exact identity relays give rooted
formula size at most
\begin{equation}
 2S+2+2(H-h-1)=2S+2(H-h).
 \label{eq:two-channel-formula-size}
\end{equation}
\end{lemma}
\begin{proof}
Add the two hidden neurons
$q_+=\operatorname{ReLU}(z_g)$ and
$q_-=\operatorname{ReLU}(-z_g)$ and use readout $q_+-q_-$.
Each new neuron has one copy of the teacher's readout formula as its
input, proving the occurrence count. Its incoming coefficients are
just the teacher's readout coefficients with a possible sign change.
Both outputs are nonnegative, so one-predecessor ReLU identity relays
preserve them exactly and add two formula occurrences per layer.
If $m=1$, the last teacher feature is nonnegative and $z_g=v h^h$ has constant sign. One added channel $\operatorname{ReLU}(|v|h^h)$ and final coefficient $\operatorname{sgn}(v)$ suffice, using only width one and no more formula occurrences. The final observation link, including clipping to $[-1,1]$, is unchanged.
\end{proof}

\begin{corollary}[Padding cost depends on scalar output, not hidden width]
\label{cor:two-channel-prior-padding}
For the fixed small-intensity L\'evy fan-in prior, let $H\ge h$ and
$N\ge m$. Define $S_*=S+H-h+1$ and
\begin{align}
 \mathcal J^{\mathrm{scalar}}_n
 =A(d+1)S_*\bigl[&(h+2)\log\{A(m+1)(B+8)\}\notag\\
 &+\log\{A S_*(d+1)(H+1)(B+8)(n+1)\}\bigr].
 \label{eq:scalar-padding-resource}
\end{align}
For conditionally sub-Gaussian labels as in the all-data expectation
oracle, the ordinary final posterior mean obeys
\begin{equation}
 \E R(\bar f_{S_n})
 \le R(g)+\frac{A(1+\sigma^2)(1+\mathcal J^{\mathrm{scalar}}_n)}{n+1}.
 \label{eq:scalar-padding-oracle}
\end{equation}
Here $g$ is the teacher's bounded observed output. The same fixed
prior and likelihood are used at every comparison resource.
\end{corollary}
\begin{proof}
When $H=h$, apply the original teacher comparison.
Otherwise apply Lemma~\ref{lem:two-channel-compiler} and then the
stable-padding comparison proof. Its relay argument tracks the
selected last-layer features and their native bulk background. The
count of new designated formula occurrences is the number of
relevant last-layer features times the number of relay layers.
There are now two such features, not $m$, so the actual count is
\eqref{eq:two-channel-formula-size}. The first-moment stability
recurrence is unchanged: each selected relay has one designated
coefficient equal to one, and the native background contracts because
its total first absolute moment is at most $1/4$.
Thus substitute teacher depth $h+1$, actual formula size $2S+2$,
and two terminal channels into the same proof. Its logarithmic
precision factor is bounded by the brackets in
\eqref{eq:scalar-padding-resource}. Applying the sharp all-data
expectation theorem gives \eqref{eq:scalar-padding-oracle}.
\end{proof}

\begin{corollary}[An additive width-and-depth resource in one dimension]
\label{cor:univariate-additive-depth}
Under the assumptions and the fixed posterior of
Theorem~\ref{thm:univariate-levy-oracle}, its conclusion improves to
\begin{equation}
 \E R(\bar f_{S_n})
 \le\inf_{\substack{C\ge1:\,6C+6\le N}}
 \left\{\inf_{g\in\mathcal G_C}R(g)
 +\frac{A(1+\sigma^2)(C+H+1)}{\alpha(n+1)}
       \log\!\bigl[e\kappa(C+1)(H+1)(n+1)\bigr]\right\}.
 \label{eq:univariate-additive-depth}
\end{equation}
\end{corollary}
\begin{proof}
The bounded-slope spline compiler has depth one, width
$m\le6C+6$, formula size $S\le m$, and
$\log(B+8)\le A\alpha^{-1}
\log[e\kappa(C+1)(n+1)]$ at $L^1$ accuracy $(n+1)^{-1}$.
Apply Corollary~\ref{cor:two-channel-prior-padding} with $h=1$ and
absorb its $4/(n+1)$ approximation error as before. If $H=1$ the
original one-layer comparison gives the same bound directly.
\end{proof}

